\chapter{Global Registration}\label{ch:global}

ICP needs a good initial alignment, which is often not available: scans taken by hand, objects
placed arbitrarily in a scene, or fragments of a broken artefact can be related by any rotation.
\emph{Global registration} finds the alignment without assuming that corresponding points are
close. Since proximity can no longer be used to find correspondences, we need more information
than the raw coordinates. Two strategies are covered, following Tema 6 and the companion notebook:
feature-based registration, which describes points by their local geometry and matches the
descriptions, and featureless registration with 4-point congruent sets (4PCS), a RANSAC-type
search that exploits affine invariants. The chapter ends with RANSAC for fitting geometric
primitives, which shares the same randomised machinery.

\section{The feature-based pipeline}

\begin{figure}[tbp]
\centering
\begin{tikzpicture}[node distance=5mm, every node/.style={font=\small},
  box/.style={draw, rounded corners=2pt, fill=MidnightBlue!8, minimum height=9mm, text width=24mm, align=center},
  ar/.style={-{Stealth[length=2.2mm]}, thick}]
\node[box] (s) {point sampling\\ \footnotesize (Harris 3D)};
\node[box, right=of s] (d) {point description\\ \footnotesize (spin images)};
\node[box, right=of d] (m) {point matching\\ \footnotesize (similarity, filtering)};
\node[box, right=of m] (t) {find transformation\\ \footnotesize (grouping, ICP)};
\draw[ar] (s) -- (d); \draw[ar] (d) -- (m); \draw[ar] (m) -- (t);
\end{tikzpicture}
\caption{The global registration pipeline based on local features.}
\label{fig:gr-pipeline}
\end{figure}

The pipeline of \cref{fig:gr-pipeline} has four stages: select a small set of distinctive points
on each shape, describe the neighbourhood of each with a descriptor invariant to rigid motions,
match descriptors between the shapes, and estimate the transformation from the matches while
discarding wrong ones. The final alignment is refined with ICP.

\section{Interest points: Harris 3D}

In images, the Harris detector \citep{Harris1988} finds corners by measuring how much a patch
changes when shifted. For an image $I$, the weighted auto-correlation of a shift $(\Delta x,\Delta
y)$ is, to first order,
\begin{equation}
c(\Delta x,\Delta y)=\sum_{(x,y)}w(x,y)\bigl(I(x+\Delta x,y+\Delta y)-I(x,y)\bigr)^2\approx
(\Delta x,\Delta y)\,\mE\,(\Delta x,\Delta y)\T ,
\end{equation}
\begin{equation}
\mE=\sum_{(x,y)}w(x,y)\begin{pmatrix}I_x^2&I_xI_y\\I_xI_y&I_y^2\end{pmatrix},
\end{equation}
with a Gaussian window $w$. The local structure is determined by the eigenvalues of $\mE$: both
small in flat regions, one large along an edge, both large at a corner. The Harris operator
$h=\det\mE-k(\tr\mE)^2$ (with $k\approx0.04$) is large only when both eigenvalues are large.

Extending this to surfaces raises two questions: how to define neighbourhoods and how to compute
derivatives. Harris 3D \citep{Sipiran2011} answers them as follows, for every vertex $v$:
\begin{enumerate}
  \item \textbf{Neighbourhood}: collect the vertices within $k$ rings (or within a radius that
  adapts to the local sampling density).
  \item \textbf{Normalisation}: translate the neighbourhood so that its centroid is at the origin and
  rotate it with principal component analysis, so that the eigenvector of smallest variance (the
  normal direction) becomes the $z$ axis; translate again so that $v$ is the origin.
  \item \textbf{Fitting}: fit a quadratic surface
  $z=f(x,y)=\tfrac{p_1}{2}x^2+p_2xy+\tfrac{p_3}{2}y^2+p_4x+p_5y+p_6$ by least squares. The surface
  now plays the role of an image.
  \item \textbf{Harris matrix}: the derivatives of $f$ are linear,
  $f_x=p_1x+p_2y+p_4$ and $f_y=p_2x+p_3y+p_5$, so the Gaussian-weighted integrals that form $\mE$
  have a closed form. With a normalised Gaussian of variance $\sigma^2$ centred at $v$,
  \begin{equation}
  \mE=\begin{pmatrix}A&C\\C&B\end{pmatrix},\quad
  \begin{aligned}
  A&=p_4^2+\sigma^2(p_1^2+p_2^2),\\ B&=p_5^2+\sigma^2(p_2^2+p_3^2),\\ C&=p_4p_5+\sigma^2(p_1p_2+p_2p_3).
  \end{aligned}
  \end{equation}
  The integration smooths the local changes of the geometry, which makes the detector robust to
  noise and tessellation.
  \item \textbf{Response and selection}: compute $h(v)=\det\mE-k(\tr\mE)^2$, keep the vertices that
  are local maxima of $h$ in their neighbourhood, and select either a fixed fraction with the
  highest response or a set spread over the surface by clustering (non-maximum suppression within a
  radius).
\end{enumerate}
The method works on meshes and point clouds. \Cref{fig:harris} shows the detected points.

\fig[0.85\textwidth]{ch08_harris.png}{Harris 3D response (red positive, blue negative) and the 60
selected interest points (yellow), with non-maximum suppression within $4\%$ of the bounding-box
diagonal. Interest points concentrate at extremities, fingers and sharp features, which are
distinctive and repeatable across scans.}{fig:harris}

Interest points are differential entities, so noise and outliers degrade them and produce wrong
correspondences. Descriptors based on integrals (integral invariants, histograms) are more robust,
and when noise, outliers and holes are all severe, featureless methods are preferable.

\section{Descriptors: spin images}

A local descriptor summarises the neighbourhood of a point in a vector that does not change under
rigid motions, so that corresponding points on two shapes have similar descriptors.
\emph{Spin images} \citep{Johnson1999} describe the distribution of surface points around an
oriented point $(\bp,\bn)$. Any point $\bx$ is expressed in the cylindrical coordinates of the local
basis formed by $\bp$, the normal $\bn$ and the tangent plane:
\begin{equation}
\alpha=\sqrt{\norm{\bx-\bp}^2-\bigl(\bn\cdot(\bx-\bp)\bigr)^2},\qquad \beta=\bn\cdot(\bx-\bp),
\end{equation}
the distance to the normal line and the signed height above the tangent plane. The angle around the
normal is discarded, which makes the pair $(\alpha,\beta)$ invariant to rotations about $\bn$:
imagine a sheet that spins around the normal and accumulates the points it sweeps. The spin image
is the 2D histogram of $(\alpha,\beta)$ over the neighbouring points, with a bin size comparable to
the sampling resolution and a support limited in $\alpha$, $\beta$, and in the angle between normals
to reduce clutter (\cref{fig:spin}). Other popular descriptors are fast point feature histograms
(FPFH) \citep{Rusu2009}, which histogram angles between normals of neighbouring pairs, and learned
descriptors.

\fig[0.85\textwidth]{ch08_spin.png}{Spin images of three points of a bunny scan (ear, neck and
flank). Each image is the histogram of the $(\alpha,\beta)$ coordinates of the neighbouring points,
with $\alpha$ horizontal and $\beta$ vertical; darker bins contain more points. Different local
shapes give clearly different images.}{fig:spin}

\section{Matching and filtering}

Two spin images $P$ and $Q$ are compared by the linear correlation coefficient over the $N$ bins
where both have data,
\begin{equation}
R(P,Q)=\frac{N\sum p_iq_i-\sum p_i\sum q_i}{\sqrt{\bigl(N\sum p_i^2-(\sum p_i)^2\bigr)\bigl(N\sum q_i^2-(\sum q_i)^2\bigr)}} .
\end{equation}
A correlation computed from few bins is unreliable, so the similarity accounts for its variance,
\begin{equation}
S(P,Q)=\bigl(\atanh R(P,Q)\bigr)^2-\lambda\frac{1}{N-3},
\end{equation}
where $\atanh$ makes the distribution of $R$ approximately normal with variance $1/(N-3)$, and
$\lambda$ weighs the penalty for small overlaps. This avoids a dependence on the density of the
data. For each sampled point of the scene, the most similar points of the model are retrieved, and
only correspondences with the highest similarity are kept.

Many correspondences are still wrong, so they are filtered in two steps:
\begin{enumerate}
  \item discard correspondences whose similarity is less than half of the maximum similarity;
  \item discard correspondences that are not geometrically consistent with enough others.
\end{enumerate}
For two correspondences $c_1=(\bs_1,\mathbf{m}_1)$ and $c_2=(\bs_2,\mathbf{m}_2)$ (scene and model points with
normals), let $S_{\mathbf{m}_2}(\mathbf{m}_1)$ denote the spin coordinates $(\alpha,\beta)$ of $\mathbf{m}_1$ in the
basis of the oriented point $\mathbf{m}_2$. A rigid motion preserves these coordinates, so consistent
correspondences satisfy $S_{\mathbf{m}_2}(\mathbf{m}_1)\approx S_{\bs_2}(\bs_1)$. The geometric consistency
distance
\begin{equation}
d_{gc}(c_1,c_2)=\frac{\norm{S_{\mathbf{m}_2}(\mathbf{m}_1)-S_{\bs_2}(\bs_1)}}{\tfrac12\bigl(\norm{S_{\mathbf{m}_2}(\mathbf{m}_1)}+\norm{S_{\bs_2}(\bs_1)}\bigr)},\qquad
D_{gc}(c_1,c_2)=\max\bigl(d_{gc}(c_1,c_2),d_{gc}(c_2,c_1)\bigr)
\end{equation}
involves both positions and normals and is small when the two correspondences are consistent.
Correspondences that are consistent ($D_{gc}$ below a threshold) with fewer than a quarter of the
list are discarded.

\section{Finding the transformation}

The surviving correspondences are grouped into mutually consistent sets. A weighted measure
favours correspondences that are far apart (they constrain the transformation better):
\begin{equation}
w_{gc}(c_1,c_2)=\frac{d_{gc}(c_1,c_2)}{1-\exp\!\bigl(-(\norm{S_{\mathbf{m}_2}(\mathbf{m}_1)}+\norm{S_{\bs_2}(\bs_1)})/(2\gamma)\bigr)},
\end{equation}
\begin{equation}
W_{gc}(c,G)=\max_{c'\in G}\max\bigl(w_{gc}(c,c'),w_{gc}(c',c)\bigr),
\end{equation}
where $W_{gc}(c,G)$ measures the consistency between a correspondence and a group $G$. For each
correspondence a group is initialised with it; the correspondence that minimises $W_{gc}(c,G)$ is
added while its value stays below a threshold, and the process continues until no more
correspondences can be added. For each group, the rigid transformation that minimises
$\sum\norm{\mR\bs_i+\bt-\mathbf{m}_i}^2$ is computed in closed form (\cref{ch:icp}), the transformations
are verified by the fraction of scene points that land on the model, and the best one is refined
with ICP.

\section{Featureless registration: 4PCS}

\subsection{From triplets to coplanar 4-point bases}

A key observation is that a pair of corresponding triplets of non-collinear points suffices to
define a unique rigid transformation. A RANSAC approach therefore picks a triplet in $P$, tries
every triplet of $Q$ as its correspondent, computes the transformation and checks how many points it
aligns. With $n$ points there are $O(n^3)$ candidate triplets per base, which is prohibitive.
\Citet{Aiger2008} showed that \emph{coplanar 4-point bases} can be matched much faster, in about
$O(n^2)$, because they have affine invariants that can be indexed.

\subsection{Measuring an alignment: the LCP}

Without correspondences, the quality of a transformation $T$ is the \emph{largest common
pointset} measure: the fraction of points of $P$ that, after transformation, have a point of $Q$
within distance $\delta$,
\begin{equation}
\mathrm{LCP}(T)=\frac{1}{\abs{P}}\,\Bigl|\bigl\{\bp\in P:\min_{\bq\in Q}\norm{T(\bp)-\bq}\le\delta\bigr\}\Bigr| .
\end{equation}
It is robust to outliers and partial overlap, because it counts agreements instead of summing
errors, and it is evaluated with a kd-tree on $Q$ (often first on a random subset, then on all the
points for the best candidates).

\subsection{Affine invariants of a coplanar base}

Let $B=\{\ba,\bb,\bc,\bd\}$ be four coplanar points such that the segments $\ba\bb$ and $\bc\bd$
intersect at a point $\be$ (\cref{fig:4pcs-inv}). The ratios
\begin{equation}
r_1=\frac{\norm{\ba-\be}}{\norm{\ba-\bb}},\qquad r_2=\frac{\norm{\bc-\be}}{\norm{\bc-\bd}}
\end{equation}
are preserved by every affine transformation, and in particular by rigid motions. Together with the
lengths $d_1=\norm{\ba-\bb}$ and $d_2=\norm{\bc-\bd}$ (preserved by rigid motions) they
characterise the base. Real data is never exactly coplanar, so the base is chosen
\emph{approximately} coplanar, with the intersection taken at the closest points of the two
lines.

\subsection{Finding congruent sets with intermediate points}

To find all 4-point sets of $Q$ congruent to $B$ (\texttt{FindCongruent}):
\begin{enumerate}
  \item collect the pairs $(\bq_i,\bq_j)$ of $Q$ whose distance is $d_1\pm\varepsilon$, and the pairs
  whose distance is $d_2\pm\varepsilon$ (both orderings of each pair);
  \item for every pair of the first set compute the \emph{intermediate point}
  $\be_1=\bq_i+r_1(\bq_j-\bq_i)$, and for every pair of the second set
  $\be_2=\bq_i+r_2(\bq_j-\bq_i)$;
  \item a pair of the first set and a pair of the second set whose intermediate points coincide,
  $\be_1\approx\be_2$, form a 4-point set with the same invariants as $B$; such coincidences are
  found with a kd-tree range search on the $\be_1$ points.
\end{enumerate}
When $d_1\approx d_2$ the two sets of pairs are the same, and both assignments of the invariants
must be tried. The candidate sets are then filtered by the angle between the segments and,
when normals are available, by the angles between normals, which are cheap and remove most false
candidates.

\fig[0.85\textwidth]{ch08_4pcs_invariants.pdf}{4PCS invariants. Left: a coplanar base whose
diagonals intersect at $\be$, with ratios $r_1=0.54$ and $r_2=0.48$. Right: in the target, the
intermediate points $\be_1$ of the pairs at distance $d_1$ (blue) and $\be_2$ of the pairs at
distance $d_2$ (orange); a coincidence of a blue and an orange point (red) reveals a congruent
4-point set.}{fig:4pcs-inv}

\subsection{The 4PCS loop}

\Cref{alg:4pcs} summarises the method. Each iteration selects a base in $P$, finds its congruent
sets in $Q$, and scores the induced transformations with the LCP. The base selection
(\texttt{SelectCoplanarBase}) looks like a detail but decides the success of the method: the base
must be \emph{wide}, because the rotation error caused by a perturbation $\varepsilon$ of the points
behaves like $\Delta\theta\sim\varepsilon/D$ for a base of diameter $D$ (\cref{fig:4pcs-base}),
yet it must lie in the overlap of the two shapes, which shrinks as the base grows. If a fraction $f$
of the points of $P$ lies in the overlap, a random base of $N$ points falls entirely in it with
probability $f^N$, so the number of iterations needed to succeed with probability $p_s$ is
\begin{equation}
L=\frac{\log(1-p_s)}{\log(1-f^N)} .
\end{equation}
4PCS returns a coarse pose; ICP finishes it (\cref{fig:4pcs-result}).

\begin{algorithm}[tbp]
\caption{4PCS global registration \citep{Aiger2008}.}\label{alg:4pcs}
\Input{point sets $P$, $Q$; overlap estimate $f$; tolerance $\delta$}
\Output{rigid transformation $T^\star$}
$T^\star\gets\mI$; $s^\star\gets0$\;
\For{$L$ iterations}{
  $B\gets$ \texttt{SelectCoplanarBase}$(P)$ (wide, approximately coplanar)\;
  compute $d_1,d_2,r_1,r_2$ of $B$\;
  $\mathcal U\gets$ \texttt{FindCongruent}$(Q,d_1,d_2,r_1,r_2)$\;
  \ForEach{$U\in\mathcal U$}{
    $T\gets$ best rigid fit of $B$ to $U$ (\cref{ch:icp})\;
    \lIf{$\mathrm{LCP}(T)>s^\star$}{$T^\star\gets T$; $s^\star\gets\mathrm{LCP}(T)$}
  }
}
refine $T^\star$ with ICP\;
\end{algorithm}

\fig[0.85\textwidth]{ch08_base_iterations.pdf}{Left: median rotation error of the transformation
estimated from a perturbed base, against the base diameter; the error decreases as $1/D$. Right:
number of iterations $L$ needed to draw at least one base inside the overlap with probability
$p_s=0.99$, for bases of $N=3$ and $N=4$ points.}{fig:4pcs-base}

\paragraph{Super4PCS.} The bottleneck of 4PCS is the extraction of all pairs at a given distance,
quadratic in the number of points. Super4PCS \citep{Mellado2014} finds them with a rasterisation of
spheres of radius $d$ around the points in a grid, and indexes the pairs by their direction to
discard incompatible angles early. The complexity becomes linear in the number of points plus the
number of candidate sets, which makes the method practical for large scans.

\fig[0.85\textwidth]{ch08_4pcs_result.png}{Global registration of two partial bunny scans in
arbitrary relative pose. Super4PCS finds an alignment with $\mathrm{LCP}=0.81$ and a rotation error
of $1.8^\circ$; ICP refines it to $0.13^\circ$.}{fig:4pcs-result}

The notebook discusses the regime for which 4PCS was designed: noisy scans with moderate overlap
and no reliable features. Its limitations, as stated by the authors, are the dependence on the
overlap estimate and the tolerance $\delta$, and degenerate geometry (planar or symmetric scenes)
where many wrong transformations obtain a high LCP.

\section{RANSAC for primitive fitting}

The random sample consensus paradigm \citep{Fischler1981} behind 4PCS also solves a different
problem: decomposing a point cloud of a man-made scene into geometric primitives. Each primitive
type has a \emph{minimal set} of points (and normals) that determines it:
\begin{itemize}
  \item \textbf{plane}: three points, or one point and its normal;
  \item \textbf{sphere}: two points with normals, since the centre lies on both normal lines (the
  midpoint of their closest points), and the radius is the mean distance to the centre;
  \item \textbf{cylinder}: two points with normals; the axis direction is
  $\bn_1\times\bn_2$, the axis passes through the intersection of the two normal lines projected on
  the plane orthogonal to the axis, and the radius is the distance of the points to the axis.
\end{itemize}
The engine repeats: draw a minimal set, fit a candidate, count its inliers (points within
distance $\varepsilon$ whose normals deviate less than $\alpha$ from the primitive's normal), and
keep the best. The iteration budget follows the same formula as for 4PCS, with $f$ the inlier
fraction of the primitive and $N$ the size of the minimal set. Safeguards against degenerate
samples (nearly parallel normals, collinear points) are essential. Several primitives are
extracted sequentially: detect the best one, remove its inliers, and repeat while the best
primitive has enough support (\cref{fig:ransac}). Efficient RANSAC \citep{Schnabel2007} samples
locally in an octree and estimates scores on subsets, and GlobFit \citep{Li2011} then enforces
global relations such as parallelism, orthogonality and coaxiality among the detected primitives.

\fig[0.7\textwidth]{ch08_ransac.png}{Sequential RANSAC on a synthetic scene of 6000 points with
$5\%$ outliers: a plane, a cylinder and a sphere are detected, in that order, and their inliers
coloured; the remaining points are outliers.}{fig:ransac}

\begin{notebook}{Tema06\_4PCS\_RANSAC.ipynb}
Part I implements 4PCS step by step: the LCP measure, the stability of wide bases
($\Delta\theta\sim\varepsilon/D$), the affine invariants, \texttt{find\_congruent} with intermediate
points, the iteration count and the main loop with two-stage verification, ICP refinement, and the
normal-angle filter. Part II applies RANSAC to planes, spheres and cylinders, with degeneracy
safeguards and sequential extraction.
\end{notebook}

\begin{keypoints}
\begin{itemize}[leftmargin=*]
  \item Global registration replaces proximity by descriptors (features) or by invariants
  (4PCS) to find correspondences.
  \item Harris 3D fits a quadric in a PCA frame and evaluates the Harris operator in closed form;
  spin images histogram $(\alpha,\beta)$ around an oriented point.
  \item Matches are filtered by similarity and geometric consistency, grouped, fitted in closed form
  and refined with ICP.
  \item 4PCS matches coplanar 4-point bases through the affine invariants $r_1,r_2$ and
  intermediate points; wide bases are stable but must fit in the overlap,
  $L=\log(1-p_s)/\log(1-f^N)$.
  \item RANSAC fits primitives from minimal sets and extracts them sequentially.
\end{itemize}
\end{keypoints}

\section*{Exercises}
\begin{exercise}
Show that the ratio $r_1$ of three collinear points is invariant under any affine map
$\bx\mapsto\mA\bx+\bb$. Why is a pair of distances not enough to characterise the base?
\end{exercise}
\begin{exercise}
How many iterations does RANSAC need to find a plane covering $20\%$ of the points with
probability $0.99$, using three-point samples? And using one point with its normal?
\end{exercise}
\begin{exercise}
Prove that the spin coordinates $(\alpha,\beta)$ of a point are invariant under rigid motions
applied to both the point and the oriented basis.
\end{exercise}
